\section{Code examples}

In this section, we will see some code examples that demonstrate how Axi works in practice. Note that Axi is not implemented yet, so the code examples are only an approximation of what we intend to achieve, rather than a demonstration of an existing language.

\begin{example}{A simple program --- list functions}

\lstinputlisting{listings/FunProg.axi}

\end{example}

The above example demonstrates the basics of Axi programming. On lines 1-3 we define an algebraic data type called \code{List}. The type is parameterized by a type \code{A} and has two constructors: \code{nil} is the empty list and \code{cons h t} is the list with head \code{h} and tail \code{t}. Both \code{:\ Type} annotations on line 1 are necessary, because we need to indicate that both \code{A} and \code{List A} are unrestricted types, as opposed to linear, affine or relevant types (we will see these in later examples).

On lines 5-8 we define a function called \code{app} which concatenates two lists. It takes three arguments: a type \code{A} and two lists \code{l1} and \code{l2} of type \code{List A}, and returns a result of type \code{List A}. The type argument means that the function is polymorphic. It is also marked with the symbol \code{@}, which means it is an \textit{implicit} argument, i.e. when calling the function it can be omitted and will get automatically inferred. The body of the function is simple: we pattern match on the argument \code{l1}. If it's \code{nil}, we return \code{l2}. If it's of the form \code{cons h t}, we return a list whose head is \code{h} and whose tail is computed by the recursive call \code{app t l2}.

On lines 10-12 we define another function, called \code{map}, which applies a function \code{f} to every element of a list. This time, however, we use an abbreviated syntax. The two type arguments \code{@A} and \code{@B} are marked as implicit, but we don't need to write the \code{:\ Type} annotations, which get inferred automatically. The argument \code{f} and its type are written out in full. Then, instead of taking another argument, we say the function returns a function of type \code{List A -> List B}. Because we omit the equality sign, what follows in the next two lines is a definition by pattern matching on the output function's \code{List A} argument.

\begin{example}{Quantitative typing}

\lstinputlisting{listings/CopyDrop.axi}

\end{example}

The above example demonstrates how the type system enforces resource discipline. We will define the type \code{KhalaniCoin} and the function \code{coinMerge} in the next example. For now it only matters that \code{KhalaniCoin} is a linear type, i.e. \code{KhalaniCoin\ :\ Type1}.

First, on line 3, we try to define a function called \code{coinForgery}. It takes an argument of type \code{KhalaniCoin} called \code{x} and returns a \code{KhalaniCoin} by calling the function \code{coinMerge} with both its arguments being \code{x}. Now, \code{KhalaniCoin} is a linear type, i.e. by default variables of this type must be used exactly once. Because \code{x} is of type \code{KhalaniCoin} and nothing indicates otherwise, this means that the function \code{coinForgery} must use \code{x} exactly once. However, it uses it twice when calling \code{coinMerge}, which is a type error.

This error is caught by the type checker, which complains that \code{x} was used too many times. The error message also hints at some possible solutions: use \code{x} less (the only solution for ordinary programmers) or redefine the type \code{KhalaniCoin} to allow duplicate uses (this solution is only available to the type's implementer). The fact that the type system was able to catch this error is crucial, because if the function \code{coinForgery} passed the type checker, it would allow programmers to create unbounded amounts of \code{KhalaniCoin}, causing hyperinflation, which most likely was not the intention of this type's creator.

The second function, called \code{coinBlackhole}, takes an argument \code{x} of type \code{KhalaniCoin} and returns \code{()}, which is the only value of type \code{Unit}, which is the type with only one element. The fact that \code{Unit} is the return type indicates that the function's result ``doesn't matter''. Because the type of \code{x} is linear, \code{coinBlackhole} must use \code{x} exactly once. However, it doesn't use it at all, which is a type error.

This error is caught by the type checker, which complains that \code{x} was not used. Similarly to the previous error message, two hints are offered: use \code{x} or make the type \code{KhalaniCoin} affine so that it allows zero or one usages (with the latter solution being available only to the type's creator). The fact that this error was caught by the type system is very important, because the creator of \code{KhalaniCoin} did not intend for values of this type to be deletable.

\begin{example}{\code{Coin} trait}

\lstinputlisting{listings/Trait.axi}

\end{example}

The above example shows the definition of the \code{Coin} trait, which could be one way of capturing the commonalities of user-defined coins. Then we implement the type \code{KhalaniCoin} from the previous example and demonstrate how to make it an instance of the \code{Coin} trait.

On lines 1-5 we define the trait \code{Coin}. It takes a parameter \code{A :\ Type1}, which is a linear type that represents the type of coins. The trait has one constant and three functions. \code{coinZero} represents a coin with denomination zero (but we will prefer the word ``amount'' instead of denomination). \code{coinMerge} merges two coins. \code{coinSplit} splits one coin into two, with its first argument being the amount of its first output and the second output being the change; there is also the possibility of failure if the amount asked is too large, indicated by \code{Option} wrapping the first output's type. \code{coinAmount} is perhaps the most mysterious --- it returns the amount of the coin, but also a new coin. This is a pattern that we will see many times. If it returned just the amount, calling \code{coinAmount} would destroy the coin, because it must be used exactly once (because its type is linear), and checking the amount counts as a single use! Therefore, \code{coinAmount} must not only return the amount, but also a ``new'' coin, which we expect to be the same as the ``old'' coin (but this property is not guaranteed; later we will see how to guarantee it using logic).

On lines 7-8 we implement the type \code{KhalaniCoin}. It is a record type with a single field \code{amount} which is a natural number (i.e. a non-negative number which can be arbitrarily large). \code{KhalaniCoin} is a linear type, because we marked it with the annotation \code{:\ Type1}. Note that this is the case even though its only field, \code{amount}, is of type \code{Nat}, which is not a linear type. In other words, the programmer is free to turn data into resources if he so wishes (although he may not freely turn resources into data).

On lines 10-29 we implement the trait \code{Coin} for the type \code{KhalaniCoin}. Note that we don't need to implement the trait's constants and functions in the same order they appear in the trait definition. On lines 11-12 we implement \code{coinZero} --- it is a record (of type \code{KhalaniCoin}) whose \code{amount} field is \code{0}.

On lines 14-15 we implement \code{coinMerge}. It takes two arguments of type \code{KhalaniCoin} and returns a \code{KhalaniCoin} which is a record whose \code{amount} field is the sum of \code{amount}s of its arguments.

On lines 17-19 we implement \code{coinAmount}. It takes an argument \code{c} of type \code{KhalaniCoin} and returns a tuple whose first component is a natural number and whose second component is a \code{KhalaniCoin}. First, on line 18, we consume the coin \code{c} by peeking at its \code{amount} field, which then gets named \code{n}. On line 19, we build the return tuple --- the first component is just \code{n}, whereas the second component is a record whose \code{amount} field is \code{n}. Note that the latter is exactly the same as the initial coin \code{c}, but we cannot directly return \code{c} as the second component, because \code{c} must be used exactly once and we have already used it when checking its \code{amount}.

Last but not least, on lines 21-29 we implement \code{coinSplit}. It takes two arguments, \code{askedAmount} of type \code{Nat} and \code{c} of type \code{KhalaniCoin}. First, on line 23, we consume \code{c} to check its amount and store it in \code{n}. Then, on line 24, we check if \code{n} is large enough for the split to succeed. If it does, on lines 25-27 we return a tuple whose first component is a coin whose amount is the \code{askedAmount} (wrapped in \code{some} to indicate success) and whose second component is a coin whose amount is \code{n} minus the \code{askedAmount} that went to the other coin. If \code{n} is not large enough, we return a tuple whose first component is \code{none} (to indicate failure) and whose second component is a coin whose amount is \code{n}, i.e. a coin that is the same as the original coin \code{c}.

\begin{example}{A simple proof}

\lstinputlisting{listings/SimpleProof.axi}

\end{example}

The above example shows a simple theorem and its proof. On line 1 we see the theorem's name, \code{app-assoc}, which is a mnemonic for ``(list) append is associative''. Its statement, on lines 2-3, reads as follows: for all types \code{A} and for all \code{l1}, \code{l2}, \code{l3} which are of type \code{List\ A}, the list \code{app\ (app\ l1\ l2)\ l3} is equal to the list \code{app\ l1\ (app\ l2\ l3)}. Note that here the symbol \code{@} marks the three arguments as \textit{explicit}, whereas its absence for \code{A} marks this type argument as implicit --- opposite to the convention used in the programming layer. These details, however, won't be very important for us now.

The proof is on lines 4-17. At any point in the proof, Axi keeps track of all the \textit{goals} that need to be solved to finish the current proof and also of the \textit{context} of these goals, which is the list of all local objects and assumptions we can use in the proof. When we start the proof with the keyword \code{proof}, we start with a single goal that is the same as the theorem statement. The code on line 5 can be read as ``let \code{A} be any type and let \code{l1} be any value of type \code{List A}''. At this point the goal becomes \code{forall (l2 l3 :\ List A), app (app l1 l2) l3 === app l1 (app l2 l3)}. On line 6 we indicate that the rest of the proof proceeds by induction on the list \code{l1}. The induction has two cases, \code{nil} and \code{cons}, corresponding to the two \code{List} constructors. This means that we now have two (sub)goals, one for each case.

On line 7 we indicate the start of the (sub)proof for the first goal, i.e. the \code{nil} case. At this point our current goal is \code{forall (l2 l3 :\ List A), app (app nil l2) l3 === app nil (app l2 l3)}, because \code{l1} was replaced with \code{nil}, and after line 8 it is \code{app (app nil l2) l3 === app nil (app l2 l3)}. Note that this information is not shown in the code --- displaying the goal and context to the programmer will be handled by language tools like IDE plugins. We finish the first subgoal on line 9. \code{refl}, which is short for ``reflexivity'', indicates that this equation holds trivially, just by computing the result on both sides.

On line 10 we indicate the start of the (sub)proof for the second goal, i.e. the \code{cons} case. On the same line we bind the arguments of \code{cons}, which we name \code{h} and \code{t}. We also get hold of our induction hypothesis for \code{t}, which we name \code{IH}. After line 11, the current goal is \code{app (app (cons h t) l2) l3 === app (cons h t) (app l2 l3)}, as \code{l1} was replaced with \code{cons h t}.

On lines 12-16, we use \textit{chaining} to finish the goal. It works like this: on line 13, after the equality symbol \code{===}, we write the left-hand side of the current goal (which must be an equation), which in our case is \code{app (app (cons h t) l2) l3} (line 13). Then, on subsequent lines, we write out what this expression is equal to (again, after \code{===}) and provide justifications (i.e. proofs) for why this is the case (after the keyword \code{by}). In our case, on line 14 we indicate that our goal's left-hand side is equal to \code{cons h (app (app t l2) l3)}, but we omit the justification for this. We could have written \code{by refl}, but this is the most trivial justification possible, so Axi guesses this is what we mean when we omit a justification. Then, on line 15, we indicate that it is equal to \code{cons h (app t (app l2 l3))}, justifying it with \code{by rewrite IH}, i.e. we are making use of our induction hypothesis. This justification is crucial and cannot be omitted, because Axi cannot infer it on its own. Finally, on line 16, we indicate that our expression is equal to the right-hand side of our goal, again omitting the justification, because it is trivial. All in all, the entire chaining proves that \code{app (app (cons h t) l2) l3} is equal to \code{app (cons h t) (app l2 l3)}, which is exactly what our goal was. Thus, the proof is finished, which we indicate on line 17 with the keyword \code{qed}.

\begin{example}{\code{LawfulCoin} concept}

\lstinputlisting{listings/Concept.axi}

\end{example}

The above example shows the definition of a \textit{concept} called \code{LawfulCoin} and its implementation for \code{KhalaniCoin}. \code{LawfulCoin} takes as a parameter one linear type \code{A}, which moreover must implement the \code{Coin} trait. This is what it means for the \code{LawfulCoin} concept to \textit{extend} the \code{Coin} trait. Note this distinction: the trait and its implementation are part of the programming layer, whereas the concept and its implementation are part of the logic layer.

On lines 2-3, we define a function \code{coinAmount'} that is part of the \code{LawfulCoin} concept. This is a feature we have not mentioned before: both traits and concepts, besides declarations, can also contain (auxiliary) definitions. In our case, we need the function \code{coinAmount'} to avoid using the much more cumbersome \code{coinAmount} when stating the properties that lawful coins should satisfy. \code{coinAmount'} takes an argument of type \code{A} (which we think of as a coin, because it has the \code{Coin} trait) and returns only the first output of \code{coinAmount}, while ignoring its second output.

It might seem that \code{coinAmount'} ignoring the coin returned by \code{coinAmount} is a type error (zero uses instead of exactly one), but function definitions that belong to the logic layer (which is indicated by the keyword \code{noncomputable}) play by different rules from ordinary functions. Most importantly, they don't need to adhere to the resource usage discipline. In other words, auxiliary functions defined solely to be used to state theorems can freely delete or copy resources, even if it were a type error for an ordinary function to do so. This is not a problem from the point of view of ordinary programs --- these ``logic-only programs'' cannot be run (which is where the \code{noncomputable} keyword comes from) --- they can only be reasoned about when proving theorems, so they can never do anything harmful to on-chain assets.

Next, on lines 5-34, comes the main part of the concept's definition. On lines 5-7 we declare a \textit{law} called \code{coinAmount-spec} which specifies the behaviour of \code{coinAmount} --- the second component of its output must be equal to its input. This is the guarantee missing from the \code{Coin} trait that we were expecting to hold. Every implementation of \code{LawfulCoin} will have to provide a proof that turns this law into a theorem.

Having given a specification of \code{coinAmount}, all the remaining laws will use \code{coinAmount'} only. On lines 9-10 we declare the law \code{coinZero-spec} which says that the amount of \code{coinZero} is zero. On lines 12-16, the law \code{coinMerge-spec} says that the amount of a merged coin is equal to the sum of amounts of the individual coins.

On lines 18-24 the law \code{coinSplit-none} says that if \code{coinSplit} returns \code{(none, coin')}, then the asked amount was not smaller than \code{coin}'s amount and the coins are equal. However, this only specifies one of the two cases that may happen when calling \code{coinSplit}, so we need another law. On lines 26-34 the law \code{coinSplit-some} says that if \code{coinSplit} returns \code{(some exact, change)}, then the asked amount was smaller than \code{coin}'s amount, the amount of \code{exact} is equal to \code{askedAmount} and the amount of \code{change} is equal to \code{coin}'s amount minus the asked amount. Together, these two laws fully specify \code{coinSplit}'s behaviour: if the coin's amount is enough, the coin gets split into a coin with the amount we asked and a coin which contains the ``change''. If we asked for too much, we just get the original coin back.

\begin{example}{Implementation of \code{LawfulCoin} for \code{KhalaniCoin}}

\lstinputlisting{listings/ConceptInstance.axi}

\end{example}

The above example shows how to implement \code{LawfulCoin} for the type \code{KhalaniCoin}. We do not need to reimplement \code{coinAmount'} or anything like that --- this function is defined by the concept and we can use it freely. What we need to do is to replace every \code{law} from the concept's definition with a \code{theorem}. It so happens that we implemented \code{coinAmount}, \code{coinZero} and \code{coinMerge} correctly, and also our implementations are so simple that the proofs go through just by computation (i.e. \code{refl}). The only theorems requiring a nontrivial amount of proving would be \code{coinSplit-none} and \code{coinSplit-some}, but we omit these proofs because they are more about natural numbers than the concept of \code{LawfulCoin}.
